\section{Conclusiones}\label{sec:conclusiones}

Buscamos una variable del lugar que explicara por qué el régimen A rinde 1.44\,t/ha más que el B, para estimar el
rendimiento en zonas sin parcelas. Con tres pruebas (diferir entre regímenes, explicar el rendimiento dentro de cada
uno y sostenerse entre vecinos) descartamos el clima, la altitud, el calendario, el suelo y la ceniza volcánica, y el
SIAP confirmó que los municipios de ambos regímenes rinden lo mismo, alrededor de 2.5\,t/ha. La brecha es una
propiedad de cómo se obtuvo la etiqueta en cada municipio, no de la zona. El modelo de la entrega sigue siendo válido
para las 59 parcelas. Para estimar en zonas nuevas harán falta datos de manejo y de productor, o rendimientos medidos
con un mismo criterio en todos los municipios.
